\cgeentry{cge_map}{Assign SAM accounts to model roles}

\cgesec{Description}

\cmd{cge_map} assigns the accounts of the SAM read by \manref{cge_sam} to the
roles of the model: activities, commodities, factors, institutions, tax
accounts, margin accounts and the other accounts described below. It then
checks every cell of the SAM against the flows that the model represents and
reports, in one table, every account that is not balanced, every account with
flows but no role, every flow that the model does not represent and every
negative cell where the model needs a non-negative value. If there is any such
problem, \cmd{cge_map} stops with error 459.

Before the checks, layouts that are common in national SAMs are rewritten into
the layout of the model, and the amounts moved are displayed. The data in
memory are not changed. With \cmd{using}, the roles are read from a template
written by \manref{cge_template}.

\cgesec{Quick start}

Map the synthetic SAM, with direct, sales and import tax accounts
\begin{stcode}
cge_map, activities(act_agr act_man act_ser) commodities(com_agr com_man com_ser) ///
    factors(lab cap) households(hhd) government(gov) savinv(sav) row(row)        ///
    dirtax(tax_dir) saltax(tax_sal) imptax(tax_imp)
\end{stcode}
As above, for a SAM that records sales and direct taxes as payments to the
government account
\begin{stcode}
cge_map, activities(act_agr act_man act_ser) commodities(com_agr com_man com_ser) ///
    factors(lab cap) households(hhd) government(gov) savinv(sav) row(row)        ///
    dirtax(gov) saltax(gov) imptax(tax_imp)
\end{stcode}
Save every cell that the checks report to \cmd{problems.csv}
\begin{stcode}
cge_map, ... report(problems.csv) replace
\end{stcode}
Map the SAM from a template filled in Excel
\begin{stcode}
cge_map using roles.xlsx
\end{stcode}

\cgesec{Syntax}

Assign the roles with options

\begin{stsyntax}
\cmdb{cge_map}\texttt{,} \optarg{}{activities}{namelist}
\optarg{}{commodities}{namelist} \optarg{}{factors}{namelist}
\optarg{}{households}{namelist} \optarg{}{government}{name}
\optarg{}{savinv}{name} \optarg{}{row}{name} [\argu{options}]
\end{stsyntax}

Assign the roles from a template

\begin{stsyntax}
\cmdb{cge_map using} \argu{template} [\texttt{,} \argu{template\_options}]
\end{stsyntax}

\begin{opttable}
\optgroup{Main}
\optarg{}{activities}{namelist} & activity accounts, each paired with the commodity in the same position of \cmd{commodities()}; required\\
\optarg{}{commodities}{namelist} & commodity accounts: the paired ones first, then any without a producing activity; required\\
\optarg{}{factors}{namelist} & factor accounts; required\\
\optarg{}{households}{namelist} & household accounts; required\\
\optarg{}{government}{name} & government account; required\\
\optarg{}{savinv}{name} & savings-investment account; required\\
\optarg{}{row}{name} & rest-of-world account; required\\
\optgroup{Institutions}
\optarg{}{enterprises}{namelist} & enterprise accounts\\
\optarg{}{insurance}{name} & social health insurance fund\\
\optgroup{Taxes}
\optarg{}{dirtax}{name} & direct tax account\\
\optarg{}{saltax}{name} & sales tax account\\
\optarg{}{imptax}{name} & import tariff account\\
\optarg{}{exptax}{name} & export tax account\\
\optarg{}{acttax}{name} & activity tax account\\
\optarg{}{vattax}{name} & value added tax (VAT) account\\
\optarg{vats}{hare}{spec} & share of each commodity's sales tax that is VAT\\
\optarg{factort}{ax}{account\texttt{:}factor \ldots} & accounts of taxes on the use of factors, each with the factor it taxes\\
\optarg{factorinc}{ometax}{name} & account of taxes on factor income\\
\optgroup{Margins, inventories and the territorial correction}
\optarg{}{marginaccounts}{namelist} & margin accounts paid per unit of composite supply\\
\optarg{domestic}{margins}{namelist} & margin accounts paid per unit of domestic sales\\
\optarg{import}{margins}{namelist} & margin accounts paid per unit of imports\\
\optarg{export}{margins}{namelist} & margin accounts paid per unit of exports\\
\optarg{}{margins}{namelist} & activities whose margins the SAM nets out of their output\\
\optarg{}{stocks}{namelist} & stock-change accounts, merged into savings-investment\\
\optarg{}{territorial}{commodity [activity]} & territorial-correction commodity and, optionally, its activity\\
\optgroup{Reporting}
\optarg{}{report}{filename} & save every cell that the checks report to a CSV file\\
\opt{}{replace} & allow \cmd{report()} to overwrite an existing file\\
\end{opttable}

\begin{opttable}[\argu{template\_options}]
\optarg{vats}{hare}{spec} & share of each commodity's sales tax that is VAT\\
\optarg{}{sheet}{sheetname} & worksheet of an Excel template; default \cmd{accounts}, or else the first sheet\\
\optarg{}{report}{filename} & save every cell that the checks report to a CSV file\\
\opt{}{replace} & allow \cmd{report()} to overwrite an existing file\\
\end{opttable}

All options other than \cmd{vatshare()}, \cmd{domesticmargins()},
\cmd{importmargins()}, \cmd{exportmargins()}, \cmd{factortax()} and
\cmd{factorincometax()} must be typed in full.

\cgesec{Options}

\begin{optdesc}
\item[\optarg{}{activities}{namelist} and \optarg{}{commodities}{namelist}]
list the activity and commodity accounts. They are paired by position: the
\argu{k}th commodity is the primary commodity of the \argu{k}th activity, used
when layouts are rewritten. An activity may also produce other commodities,
in the proportions that the SAM records. Each activity needs a paired
commodity; commodities without a producing activity, such as commodities that
are only imported, are listed after the paired ones. A payment from a
household to an activity, $\SAM(\textit{activity}, \textit{household})$, is
consumption of own produce at home (\secref{sec:home}).

\item[\optarg{}{factors}{namelist} and \optarg{}{households}{namelist}] list
the factor and household accounts. Non-profit institutions serving households
can be listed as households.

\item[\optarg{}{government}{name}, \optarg{}{savinv}{name} and
\optarg{}{row}{name}] name the government, savings-investment and
rest-of-world accounts.

\item[\optarg{}{enterprises}{namelist}] lists enterprise accounts
(corporations). Enterprises receive factor income and transfers, pay direct
taxes and save; the rest of their income is transferred to other institutions
or abroad. They do not consume, so a purchase of commodities by an enterprise
is reported as a flow that the model does not represent.

\item[\optarg{}{insurance}{name}] names a social health insurance fund. The
fund may receive contributions from factors (its row, the factors' columns),
transfers from the government and from the rest of the world, and may buy
commodities and save; it may have no other flows. Contributions enter the
model as a fixed rate on each factor's income, taken before that income is
distributed to institutions, and the fund buys a fixed basket of commodities
with its revenue less a fixed share saved, as in equations \eqref{eq:I2} and
\eqref{eq:I12} to \eqref{eq:I14}. A non-profit institution financed by donors that buys health
services is mapped as a household instead.

\item[\optarg{}{dirtax}{name}, \optarg{}{saltax}{name},
\optarg{}{imptax}{name}, \optarg{}{exptax}{name} and \optarg{}{acttax}{name}]
name the accounts that collect direct taxes, sales taxes, import tariffs,
export taxes and activity taxes. The rates of a tax that is not mapped are
zero. If the SAM records a tax as a payment to the government account itself,
name the government account in the option, for example \cmd{saltax(gov)}: the
government row's cells in the columns that pay the tax (households and
enterprises for \cmd{dirtax()}, commodities for \cmd{saltax()},
\cmd{imptax()} and \cmd{exptax()}, activities for \cmd{acttax()}) are then
routed through an internal tax account, \cmd{_dirtax}, \cmd{_saltax},
\cmd{_imptax}, \cmd{_exptax} or \cmd{_acttax}. At most one of \cmd{saltax()},
\cmd{imptax()}, \cmd{exptax()} and \cmd{vattax()} may name the government
account, because their cells would otherwise be ambiguous.

\item[\optarg{}{vattax}{name}] names the account that collects VAT from the
commodities, net of input credits, and pays it to the government. VAT is then
a tax on final demand at a rate by commodity, with credits for the VAT on
inputs in the shares given by \cmd{vatcredit()} of \manref{cge_calibrate};
exports and margin services are zero-rated (\secref{sec:vat}). The government
account may be named, as for the other taxes; the VAT is then routed through
the internal account \cmd{_vattax}.

\item[\optarg{vats}{hare}{spec}] is for SAMs that record VAT within the sales
tax account. It moves the given share of each commodity's sales tax to an
internal VAT account, \cmd{_vatshare}, and the same amount of the sales tax
account's payment to the government to the payment of \cmd{_vatshare}.
\argu{spec} is \argu{commodity}\cmd{:}\argu{share} pairs, or one share for
every commodity; shares lie between 0 and 1, and commodities not listed have
share 0. \cmd{vatshare(0)} creates an empty VAT account, so that a VAT can be
introduced by a shock to \cmd{tv}. \cmd{vatshare()} needs \cmd{saltax()} and
may not be combined with \cmd{vattax()}.

\item[\optarg{factort}{ax}{account\texttt{:}factor \ldots}] names accounts of
taxes on the use of factors, such as payroll taxes on labour or taxes on
capital, which activities pay and which pay the government. Each account is
followed by the factor whose use it taxes. The rate on each factor in each
activity is the tax over the factor payment; it raises the cost of the factor
to the activity, and value added at factor cost includes it
(\secref{sec:factortax}).

\item[\optarg{factorinc}{ometax}{name}] names the account of taxes on factor
income, which factors pay and which pays the government. The rate on each
factor's income is taken before the income is distributed to institutions, as
the contributions to a health insurance fund are.

\item[\optarg{}{marginaccounts}{namelist}] lists margin accounts: accounts that
collect trade or transport margins from commodities (the margin account's row,
the commodities' columns) and spend them on margin services (the commodities'
rows, the margin account's column). A margin account may have no other flows.
Margins enter the model as a fixed quantity of margin services per unit of
each commodity's composite supply, and their cost enters the commodity's
price, as in equations \eqref{eq:P3} and \eqref{eq:T8}.

\item[\optarg{domestic}{margins}{namelist}, \optarg{import}{margins}{namelist}
and \optarg{export}{margins}{namelist}] list margin accounts whose margins are
paid per unit of domestic sales, of imports and of exports, for SAMs that
separate them (\secref{sec:margins-market}). Domestic buyers then pay the
producer price plus domestic margins, the import price includes import
margins, and exporters receive the world price less export taxes and export
margins. A commodity that pays such margins must have domestic sales, imports
or exports respectively; \manref{cge_calibrate} checks this.

\item[\optarg{}{margins}{namelist}] lists activities, typically trade and
transport, whose margin revenue the SAM nets out of their output by recording
a negative use of their own product. That negative cell is moved to the
activity's output, so that the margins enter the model as the purchases of
trade and transport services by the activities that pay them, as the SAM
records them. The activity's own use of its product cannot be separated from
the netted margins and becomes zero. An activity named here must be one of
the activities, and must have a negative use of its paired commodity.

\item[\optarg{}{stocks}{namelist}] lists stock-change accounts: accounts
financed by savings-investment that buy commodities, with negative cells for a
fall in inventories. Their purchases are moved to the savings-investment
account, so that changes in inventories become part of investment demand. A
stock-change account may receive payments only from savings-investment and may
have no other role.

\item[\optarg{}{territorial}{commodity [activity]}] names the
territorial-correction commodity and, optionally, its activity. The
correction's imports (purchases abroad by residents) become household payments
to the rest of the world, and its exports (purchases in the country by
non-residents) become transfers from the rest of the world to households, both
divided among households in proportion to their consumption of the correction.
The correction may have only household consumption, imports and exports, and
the activity only exports, which are first moved to the commodity. Both
accounts then have no flows, and may have no other role.

\item[\optarg{}{report}{filename}] saves every cell that the checks report to a
CSV file with the columns \cmd{problem}, \cmd{paid_by}, \cmd{paid_by_role},
\cmd{received_by}, \cmd{received_by_role} and \cmd{value}. The file is written
whether or not problems are found.

\item[\opt{}{replace}] allows \cmd{report()} to overwrite an existing file.

\item[\optarg{}{sheet}{sheetname}] names the worksheet of an Excel template.
The default is the sheet \cmd{accounts}, which \manref{cge_template} writes,
or else the first sheet. It may be given only with \cmd{using}.
\end{optdesc}

\cgesec{Remarks and examples}

The mapping tells the model what each account of the SAM is. Each account has
one role, given by the option that lists it, except that the government
account may also be named in tax options and a margin activity is listed in
both \cmd{activities()} and \cmd{margins()}. An account with no flows needs no
role. The mapping replaces any earlier mapping, and the model must be
calibrated again with \manref{cge_calibrate}.

\cgesubsec{The checks}

After the layouts are rewritten, \cmd{cge_map} checks the SAM and collects four
kinds of problem:
\begin{itemize}
\item an account that is not balanced: its row total differs from its column
total;
\item an account with flows but no role;
\item a non-zero cell that is not one of the flows of the model, listed below;
\item a negative cell where the model needs a non-negative value.
\end{itemize}
The problems are summarised in one table, grouped by problem and by the roles
of the paying and receiving accounts, the largest group first, with the number
of cells, their sum in SAM units and their share of the sum of the absolute
values of all cells. For an account with no role, the value is its row total,
and its flows are also listed by role. Unknown account names (error 111) and an
account given two roles (error 198) stop the command before the checks.

The model represents the following payments, from the column account to the
row account:

\begin{cgetable}{>{\raggedright\arraybackslash}p{0.26\linewidth}>{\raggedright\arraybackslash}p{0.70\linewidth}}
\toprule
Paid by (column) & Received by (row)\\
\midrule
commodity & activities (domestic output); rest of the world (imports); sales, import and export tax accounts; VAT account; margin accounts\\
activity & commodities (intermediate inputs); factors; activity tax account; accounts of taxes on factor use\\
factor & households; enterprises; government; rest of the world; health insurance fund; account of taxes on factor income\\
household & commodities; activities (home consumption); direct tax account; savings-investment; households and enterprises; rest of the world\\
enterprise & direct tax account; savings-investment; households and enterprises; rest of the world\\
government & commodities; households and enterprises; government; rest of the world; savings-investment; health insurance fund\\
tax account & government\\
savings-investment & commodities (investment, including changes in inventories)\\
margin account & commodities (margin services)\\
health insurance fund & commodities; savings-investment\\
rest of the world & commodities (exports); households, enterprises and government; factors (factor income from abroad); savings-investment; health insurance fund\\
\bottomrule
\end{cgetable}

Negative cells are allowed in the rows of the tax accounts and of
savings-investment, in transfers (the payments among households, enterprises,
the government and the rest of the world, and factor income from abroad), in
the payments of tax accounts to the government, in investment demand (a fall in
inventories) and in sales of commodities by activities. A negative sale gives
a negative yield; it is displayed as a note and counted in
\res{r(N_negative_sales)}. Section~\ref{sec:flows} gives the rules in full.

\cgesubsec{Layouts rewritten before the checks}

National SAMs often record flows differently from the model. \cmd{cge_map}
rewrites these layouts, in this order, and displays the amounts moved:

\begin{cgetable}{>{\raggedright\arraybackslash}p{0.40\linewidth}>{\raggedright\arraybackslash}p{0.17\linewidth}>{\raggedright\arraybackslash}p{0.37\linewidth}}
\toprule
The SAM records & Option & Rewrite\\
\midrule
exports as sales by activities to the rest of the world & (always) & moved to the commodity paired with the activity\\
margins netted out of margin activities, as a negative use of their own product & \cmd{margins()} & moved to the activity's output\\
a territorial correction & \cmd{territorial()} & household payments to, and transfers from, the rest of the world\\
changes in inventories in a stock-change account & \cmd{stocks()} & merged into savings-investment\\
taxes paid to the government account & tax option naming the government & routed through an internal tax account\\
VAT within the sales tax & \cmd{vatshare()} & moved to the internal account \cmd{_vatshare}\\
\bottomrule
\end{cgetable}

The rewrites leave unchanged every total that enters the model, except that a
margin rewrite raises the output of the margin activity, and the supply and
demand of its product, by the margins, and a territorial rewrite replaces the
correction's imports, exports and consumption by transfers
(\secref{sec:layouts}). \manref{cge_samexport} with \cmd{mapped} writes the SAM
after the rewrites.

\cgesubsec{Example 1: the standard layout}

The synthetic SAM has separate tax accounts and records exports as sales of
commodities:
\stlog{map_standard}

\cgesubsec{Example 2: a problem, and a layout rewrite}

The national-accounts layout of the synthetic SAM describes the same economy,
but records exports as sales by activities, and sales and direct taxes as
payments to the government account. Mapped with an import tax account only, it
gives two groups of problems:
\stlog[\scriptsize]{map_problem}
The exports of the activities, 75 in all, were moved to their paired
commodities before the checks. The cells left are payments by the commodities
and by the household to the government, which the model represents only
through tax accounts. With \cmd{report()}, every cell of the table was saved:
\stlog{map_cells}
Naming the government account in \cmd{saltax()} and \cmd{dirtax()} routes the
taxes through internal accounts, and the mapping succeeds:
\stlog{map_national}
\manref{cge_calibrate} gives the same model for this layout as for the
standard one.

\cgesubsec{Example 3: a template}

A template lists the roles in a file, one row per account, which is easier to
check and to keep than a long command. Written after a mapping, it holds that
mapping, and \cmd{cge_map using} maps the SAM from it and displays the
equivalent command:
\stlog{map_template}
The government account has three roles, \cmd{government dirtax saltax}, so the
taxes are routed through internal accounts as in example~2. Because the model
is kept in Mata, the data in memory, here the template, can be replaced at
any time.

In a template, the \cmd{role} column names the options that list the account,
separated by spaces or commas, in either case; the singular forms
\cmd{activity}, \cmd{commodity}, \cmd{factor}, \cmd{household},
\cmd{enterprise}, \cmd{margin}, \cmd{marginaccount}, \cmd{stock},
\cmd{vat}, \cmd{domesticmargin}, \cmd{importmargin} and \cmd{exportmargin} are
also accepted. The accounts of each option are listed in the order of the SAM,
and the correction commodity of \cmd{territorial} before its activity. The
\cmd{pair} column pairs an activity with a commodity; where it is blank, the
activity is paired with the commodity whose column pays it most. A tie, an
activity with no sales to commodities, and a commodity paired with two
activities are errors. For an account of taxes on factor use, the \cmd{pair}
column names the factor it taxes. The template may be a \cmd{.csv},
\cmd{.xls}, \cmd{.xlsx} or \cmd{.dta} file with a string variable
\cmd{account}.

If the SAM was aggregated with the same template (\manref{cge_aggregate}), its
accounts are the template's groups: each group takes the role of its members,
which must agree, and is paired with the group of its members' pairs.

\cgesubsec{Example 4: a supply-use layout}

The supply-use layout of the synthetic SAM has two margin accounts, a
commodity that is only imported, \cmd{com_fuel}, listed after the paired
commodities, and a stock-change account:
\stlog{map_sut}
The stock-change account received 3 from savings-investment; its purchases are
now investment demand.

\cgesec{Stored results}

\cmd{cge_map} stores the following in \cmd{r()}:
\begin{storedres}
\resgroup{Scalars}
\res{r(N_negative_sales)} & number of negative sales of commodities by activities\\
\res{r(exports_moved)} & exports recorded as sales by activities, moved to commodities\\
\res{r(margins_moved)} & margins moved to the output of margin activities\\
\res{r(territorial_imports)} & purchases abroad by residents, made household payments abroad\\
\res{r(territorial_exports)} & purchases by non-residents, made transfers from abroad to households\\
\res{r(stocks_merged)} & payments of savings-investment to stock-change accounts\\
\res{r(taxes_rerouted)} & direct, sales, import, export and activity taxes routed through internal tax accounts\\
\res{r(vat_moved)} & VAT moved out of the sales tax by \cmd{vatshare()}\\
\resgroup{Macros}
\res{r(command)} & with \cmd{using}, the equivalent \cmd{cge_map} command\\
\end{storedres}

\cgesec{Methods and formulas}

Let $\SAM(r, c)$ be the payment from column account $c$ to row account $r$
before a rewrite and $\SAM'(r, c)$ after it, $c(a)$ the commodity paired with
activity $a$, and \textit{row}, \textit{sav} and \textit{gov} the
rest-of-world, savings-investment and government accounts. Cells not shown are
unchanged. The rewrites are made in the following order.

Exports recorded as sales by activity $a$:
\begin{gather*}
\SAM'(c(a), \textit{row}) = \SAM(c(a), \textit{row}) + \SAM(a, \textit{row}), \qquad
\SAM'(a, c(a)) = \SAM(a, c(a)) + \SAM(a, \textit{row}),\\
\SAM'(a, \textit{row}) = 0 .
\end{gather*}
Margins netted out of activity $a$, with $\SAM(c(a), a) < 0$:
\begin{equation*}
\SAM'(a, c(a)) = \SAM(a, c(a)) - \SAM(c(a), a), \qquad \SAM'(c(a), a) = 0 .
\end{equation*}
Territorial correction $k$: the exports of its activity, if any, are first
moved to $k$ as above. With $\mathit{IM} = \SAM(\textit{row}, k)$,
$\mathit{EX} = \SAM(k, \textit{row})$ and the shares
$s_h = \SAM(k, h) / \sum_{h'} \SAM(k, h')$ of the households,
\begin{equation*}
\SAM'(\textit{row}, h) = \SAM(\textit{row}, h) + s_h\,\mathit{IM}, \qquad
\SAM'(h, \textit{row}) = \SAM(h, \textit{row}) + s_h\,\mathit{EX},
\end{equation*}
and the cells of $k$ are set to zero. Stock-change account $k$:
$\SAM'(r, \textit{sav}) = \SAM(r, \textit{sav}) + \SAM(r, k)$ for every account
$r$, and the row and column of $k$ are set to zero.

A tax paid to the government by the accounts $j \in P$ (the payers of that
tax) is routed through a new account $t$:
\begin{equation*}
\SAM'(t, j) = \SAM(\textit{gov}, j), \qquad \SAM'(\textit{gov}, j) = 0, \qquad
\SAM'(\textit{gov}, t) = \sum_{j \in P} \SAM(\textit{gov}, j) .
\end{equation*}
With VAT shares $v_c$ of the sales tax account \textit{sal}, the new account
$t$, \cmd{_vatshare}, has
\begin{gather*}
\SAM'(t, c) = v_c\,\SAM(\textit{sal}, c), \qquad
\SAM'(\textit{sal}, c) = (1 - v_c)\,\SAM(\textit{sal}, c),\\
\SAM'(\textit{gov}, t) = \sum_c v_c\,\SAM(\textit{sal}, c), \qquad
\SAM'(\textit{gov}, \textit{sal}) = \SAM(\textit{gov}, \textit{sal}) - \SAM'(\textit{gov}, t) .
\end{gather*}

The checks use the tolerance $\tau = 10^{-8} \max\bigl(1, \max_{r,c}
|\SAM(r, c)|\bigr)$ on the rewritten SAM. A cell is non-zero if
$|\SAM(r, c)| > \tau$, and account $i$ is not balanced if
\begin{equation*}
\Bigl|\sum_c \SAM(i, c) - \sum_r \SAM(r, i)\Bigr| > \tau .
\end{equation*}
The share of a group of problem cells $G$ in the report is
$\sum_{(r,c) \in G} |\SAM(r, c)| \big/ \sum_{r,c} |\SAM(r, c)|$.

Section~\ref{sec:flows} lists the flows of the model and the cells that may be
negative, and \secref{sec:layouts} gives the rewrites and the totals they
change.

\cgesec{Also see}

\begin{alsosee}
\seeentry{cge_sam}{Read a social accounting matrix}
\seeentry{cge_template}{Write a template that maps SAM accounts to roles and groups}
\seeentry{cge_aggregate}{Merge SAM accounts into groups}
\seeentry{cge_samexport}{Write the SAM to a file or to memory}
\seeentry{cge_calibrate}{Calibrate the model to the SAM}
\seechapter{chap:calibration}{Calibration}
\end{alsosee}
